% !PRINTED 45
\weberpage{45}
%<<CONT

but treat it now under the special supposition that the given values of $x$ are the first $n+1$ whole numbers $0,1,2,\ldots,n$.

The binomial coefficients $B_\nu^{(n)}$, as they were defined by \S~7, (6), keep their meaning even when $n$ is not a whole number, as was previously presupposed, but an arbitrary variable quantity. Then
\begin{equation}
 B_\nu^{(x)}={x(x-1)\cdots(x-\nu+1)\over 1\cdot2\cdot3\cdots\nu}
\tag{3}
\end{equation}
is a whole rational function of degree $\nu$ in $x$.
\footnote{The meaning of these generalized binomial coefficients for the expansion of powers of a binomial does not belong here, but to analysis.}
% SOURCE-ORIGINAL-SIC: The English regularizes the source footnote's misspelling "Binominalcoefficienten"; the German leaf preserves it verbatim.
The coefficient of $x^\nu$ is different from zero, and therefore $x^\nu$ can also be expressed by $B_\nu^{(x)}$ and by lower powers of $x$. Thus every whole function $f(x)$ of degree $n$ in $x$ can also be linearly expressed by $B_0^{(x)},B_1^{(x)},\ldots,B_n^{(x)}$ in the form
\begin{equation}
 f(x)=M_0B_0^{(x)}+M_1B_1^{(x)}+\cdots+M_nB_n^{(x)},
\tag{4}
\end{equation}
where $M_0,M_1,\ldots,M_n$ are constants; and the function $f(x)$ is determined when these constants are determined.

Now, by our supposition, $f(0),f(1),f(2),\ldots,f(n)$ are given. Since, by (3), $B_\nu^{(x)}$ always vanishes when $x$ has one of the values $0,1,2,\ldots,\nu-1$, the following linear equations for the unknowns $M$ are obtained from (4):
\begin{equation}
\begin{aligned}
 f(0)&=M_0B_0^{(0)},\\
 f(1)&=M_0B_0^{(1)}+M_1B_1^{(1)},\\
 f(2)&=M_0B_0^{(2)}+M_1B_1^{(2)}+M_2B_2^{(2)},\\
 &\quad\cdots\\
 f(n)&=M_0B_0^{(n)}+M_1B_1^{(n)}+M_2B_2^{(n)}+\cdots+M_nB_n^{(n)}.
\end{aligned}
\tag{5}
\end{equation}
These equations are now to be solved with respect to $M_0,M_1,\ldots,M_n$, which is done very easily with the help of the final equations (8) of the preceding paragraph. The first equation (5) gives directly
\begin{equation}
 M_0=f(0).
\tag{6}
\end{equation}
